\section{Background and Related Work}
\label{sec:background}

\textbf{Aggregated serving.}
Orca introduced iteration-level scheduling, and vLLM's PagedAttention enables
memory-efficient continuous batching~\cite{yu2022orca,kwon2023vllm}. Long
prefills can nevertheless delay generation. Sarathi-Serve bounds prefill chunks
and constructs decode-maximal mixed batches; Dynamic SplitFuse similarly splits
and combines prompt and generation work~\cite{agrawal2024sarathi,
holmes2024fastgen}. PD-TDM retains this prior bounded-prefill contribution. Its
additional control is the opportunity assigned to each phase across iterations.
Our measured baseline is the native chunked-prefill path in vLLM-Ascend on the
same NPUs, denoted \baseline{}; we do not claim a cross-platform comparison to
the Sarathi-Serve system.

\textbf{Disaggregation and spatial multiplexing.}
Splitwise and DistServe place prefill and decode on different instances, while
Mooncake builds a KV-cache-centric distributed architecture~\cite{
patel2024splitwise,zhong2024distserve,qin2025mooncake}. These designs gain
independent phase capacity at the cost of additional replicas and state
movement. MuxWise and semi-PD share storage while concurrently partitioning GPU
compute between the phases~\cite{chen2026muxwise,hong2025semipd}. Their spatial
isolation is orthogonal to PD-TDM, which gives the entire TP replica to one
phase at a time.

\textbf{Temporal orchestration.}
EcoServe is the closest prior design: it temporally disaggregates phases within
an instance and rolls activation across multiple cooperating instances to keep
prefill available~\cite{du2026ecoserve}. PD-TDM does not claim to originate
temporal disaggregation. It instead studies bounded, iteration-granular phase
scheduling when the deployment has only one shared TP replica, and explicitly
reports the decode-pressure region in which this choice fails.
